\section{Discussion}
\label{sec:discussion}

\para{What do the results mean?} The cleanest conclusion is that LayerNorm acts as a geometry-shaping component in GPT-2-class models. On \ocw, removing it produces less anisotropic and higher-rank clue embeddings for every wall in the evaluation set. This is exactly the direction expected if LayerNorm helps compress the residual stream into a more homogeneous geometry. The generation results partially agree with that story: semantic spread increases under fixed decoding, which suggests that the model explores a broader region of concept space when LayerNorm is absent.

\para{Why does broader geometry not become broader lexical novelty?} Our results suggest that representation breadth is necessary but not sufficient. The LN-free model may access more varied hidden states while still decoding into similar lexical forms because the pretrained distribution, model capacity, and weak prompt following of GPT-2 remain the dominant bottlenecks. In other words, a broader internal manifold does not automatically become better controlled creativity at the output surface.

\para{Alternative explanations.} One possibility is that \lnfreegpt is simply noisier after post-training LayerNorm removal. That could increase embedding dispersion without improving useful novelty. A second possibility is benchmark mismatch: \noveltybench was designed for stronger instruction-following models, so GPT-2 may be too weak for lexical-diversity gains to emerge cleanly. A third possibility is measurement mismatch. Our \ocw probe uses direct clue embeddings from a causal language model rather than a task-trained grouping head, so it may reveal geometry more readily than task competence.

\para{Limitations.} This study uses a single model family, a small released LN-free derivative, and only one normalization ablation. The sentence-adherence metric is based on just three prompts and is therefore too sparse to carry much weight. The geometry benchmark is indirect, because we evaluate clue embeddings rather than a supervised association model. Finally, the experiments ran on CPU after a CUDA driver mismatch, which did not affect correctness but limited the scope of possible follow-up sweeps.

\para{Broader implications.} These findings narrow the claim we can make. They do not show that LayerNorm is the main cause of creative collapse in modern aligned language models. They do show that normalization changes internal geometry in a direction consistent with representational compression, and that removing it alone does not solve the harder problem of turning latent variation into controlled, useful novelty. For creativity-oriented model design, that points toward combined interventions: architecture, decoder, and evaluation must likely move together.
